% =============================================================================
% 01_prediction_design.tex — Problem 1: Prediction Design and Leakage Control
% =============================================================================
\section{Prediction Design and Leakage Control}
\label{sec:prediction-design}

% ---- 1.1 Target and Predictors ---------------------------------------------
\subsection{Target and Predictors}
\label{sec:target-predictors}

The \texttt{fat.csv} dataset contains $n = 252$ observations and 18 variables
describing body composition and anthropometric measurements of adult males.
The response variable is \texttt{brozek}, the body-fat percentage estimated via
Brozek's equation.

After inspection, the following variables are \textbf{excluded} from the
predictor set to prevent data leakage:
%
\begin{itemize}
    \item \texttt{siri} — an alternative body-fat estimate; including it would
          constitute a near-perfect proxy for the response.
    \item \texttt{density} — body density is the direct input to both the Brozek
          and Siri formulas; using it makes prediction trivially deterministic.
    \item \texttt{free} — fat-free weight is algebraically derived from body fat
          and total weight, creating a circular dependency.
\end{itemize}
%
% TODO: Justify each exclusion in more detail. Explain the algebraic
%       relationship between brozek, siri, density, and free. Cite the
%       original Brozek / Siri formulas if possible.

This leaves $p = 14$ candidate predictors (age, weight, height, adipos, neck,
chest, abdom, hip, thigh, knee, ankle, biceps, forearm, wrist).

\paragraph{Summary statistics.}
Table~\ref{tab:summary-stats} reports the descriptive statistics for all
retained predictors computed on the \emph{training set only}.

\begin{table}[H]
\centering
\caption{Descriptive statistics of retained predictors (training set).}
\label{tab:summary-stats}
% TODO: Uncomment the line below once the table file is generated by R.
% \input{../output/tables/tab_p1_summary}
%
% Placeholder while the pipeline has not been run:
\fbox{\parbox{0.8\textwidth}{\centering\vspace{1cm}
    \texttt{tab\_summary\_stats.tex} — run the R pipeline to generate.
\vspace{1cm}}}
\end{table}

% ---- 1.2 Split and Train-Only Preprocessing --------------------------------
\subsection{Split and Train-Only Preprocessing}
\label{sec:split-preprocess}

The data are split into a training set ($n_{\mathrm{train}} \approx 201$,
$\sim 80\%$) and a held-out test set ($n_{\mathrm{test}} \approx 51$,
$\sim 20\%$) using a fixed random seed for reproducibility.

% TODO: Report the exact seed value used, e.g.:
% \texttt{set.seed(2025)}

\paragraph{Standardisation.}
All predictors are centred and scaled using the \emph{training-set} mean and
standard deviation.  The same transformation is then applied to the test set to
prevent information leakage from the test data into the modelling pipeline.
Formally, for each predictor $j$:
\[
    \tilde{x}_{ij}
    = \frac{x_{ij} - \bar{x}_j^{\mathrm{train}}}
           {s_j^{\mathrm{train}}},
    \qquad i = 1, \dots, n.
\]

\paragraph{Cross-validation.}
We use $K = 5$-fold cross-validation (CV) throughout, with folds assigned once
and reused across all models to ensure fair comparison.

% TODO: Discuss additional leakage safeguards taken in the pipeline
%       (e.g., not using test labels anywhere, scaling inside CV folds, etc.).

% ---- 1.3 Training-Data Exploration -----------------------------------------
\subsection{Training-Data Exploration}
\label{sec:eda}

% -- Correlation heatmap
\includefigure{0.85\textwidth}{fig_corr_heatmap}%
    {Pearson correlation heatmap of training-set predictors and response.}%
    {fig:corr-heatmap}

% TODO: Interpret the heatmap. Which predictor pairs show high collinearity?
%       How does this motivate regularisation?

% -- Response distribution
\includefigure{0.65\textwidth}{fig_p1_dist}%
    {Distribution of the response variable \texttt{brozek} in the training set.}%
    {fig:response-dist}

% TODO: Comment on skewness, outliers, or the need for transformation.

% -- Boxplots of predictors
\includefigure{0.85\textwidth}{fig_p1_boxplots}%
    {Side-by-side boxplots of standardised predictors (training set).}%
    {fig:boxplots}

% TODO: Identify at least one anomaly or unusual observation.

% -- Pairwise scatter
\includefigure{0.95\textwidth}{fig_p1_pairwise_scatter}%
    {Pairwise scatter plots of selected predictors against \texttt{brozek}.}%
    {fig:pairwise}

% TODO: Discuss linearity assumptions and whether any transformations are needed.
%       Summarise how multicollinearity motivates the use of Ridge / Lasso.
